\documentclass[11pt]{article}
\usepackage[margin=1in]{geometry}
\usepackage{amsmath,amssymb}
\usepackage{booktabs}
\usepackage[dvipsnames]{xcolor}
\usepackage{hyperref}
\hypersetup{colorlinks=true,linkcolor=blue!50!black,urlcolor=blue!50!black,citecolor=blue!50!black}

\title{Station Moves Matter:\\
Documented Relocations and the Frankfurt Urban--Reference\\
Temperature Gap, 1985--2025}
\author{David Lindgreen}
\date{October 9, 2026 \\[2pt] \small Version 0.1.0 \quad Preprint, not peer reviewed}

\emergencystretch=3em
\begin{document}
\maketitle

\begin{abstract}
We re-analyse the daily temperature difference between two Deutscher Wetterdienst (DWD) stations that DWD itself
designates as an urban/reference pair, Frankfurt/Main-Westend and Frankfurt/Main (14{,}579 paired days,
1985--2025). The preregistered result reproduces exactly from freshly downloaded archives: a mean daily-mean gap of
$+0.455\,^{\circ}$C (95\% block-bootstrap interval 0.432 to 0.478) and no detectable linear trend. A post-hoc check of two
limitations the original report only mentioned shows the single average is not a stable property of the pair. The
gap changed by $-0.43\,^{\circ}$C (95\% interval $-0.54$ to $-0.29$) across the urban station's 2008 relocation and by
$+0.19\,^{\circ}$C (0.10 to 0.27) across the reference station's 2014 relocation, each comparable to the whole mean
gap, and our earlier statement that neither relocation produced an unusual jump was wrong. The daily-minimum gap
($+1.12\,^{\circ}$C) is about 2.5 times the daily-mean gap and the daily-maximum gap ($+0.12\,^{\circ}$C) is barely
positive. A null trend is therefore uninformative while these steps remain unmodelled.
\end{abstract}

\section{Question}
How much warmer is urban Frankfurt than DWD's own designated reference counterpart, with what uncertainty, and has the
gap trended over four decades? We use daily mean temperature (TMK) at station 01424 (Frankfurt/Main-Westend, urban) and
station 01420 (Frankfurt/Main, an airport that DWD labels the counterpart). The gap is the urban minus the reference value on
days with both observations. This is a statement about two stations, not about the urban heat island of the whole city.

\section{Preregistered analysis}
The protocol fixed the stations, the period (1985-11-01 to 2025-12-31), the rule for missing days, a 30-day moving-block
bootstrap for the mean gap (daily gaps are strongly autocorrelated), and an ordinary least-squares trend on annual mean gaps
for the 40 years with at least 300 valid paired days. A post-release Newey--West (three-lag) interval~\cite{newey1987} was
added as a sensitivity because the annual residuals have lag-1 correlation 0.657.
\begin{table}[h]
\centering\small
\begin{tabular}{lrrl}
\toprule
Period & $n$ days & Mean gap ($^{\circ}$C) & 95\% interval\\
\midrule
Full period & 14{,}579 & +0.455 & 0.432 to 0.478\\
Winter (DJF) & 3{,}641 & +0.489 & 0.453 to 0.522\\
Spring (MAM) & 3{,}619 & +0.541 & 0.497 to 0.589\\
Summer (JJA) & 3{,}649 & +0.415 & 0.361 to 0.471\\
Autumn (SON) & 3{,}670 & +0.374 & 0.341 to 0.406\\
\bottomrule
\end{tabular}
\caption{Preregistered mean gap (daily mean temperature) with block-bootstrap intervals.}
\end{table}
The annual-mean trend is $-0.0031\,^{\circ}$C per year (classical 95\% interval $-0.0071$ to $+0.0008$, $p=0.118$; Newey--West
$-0.0079$ to $+0.0016$, $p=0.186$): no linear trend is detected. All of this regenerates identically from a fresh download
of the two DWD archives.

\section{Post-hoc extensions}
The original report named two untested limitations: the daily minimum and maximum (TNK, TXK) might show different gaps, and
the documented relocations (urban 2008-07-01, reference 2014-10-22) might introduce steps. We tested both with the same
block bootstrap; none of this was preregistered and every result is reported.
\begin{table}[h]
\centering\small
\begin{tabular}{lccc}
\toprule
Variable & Mean gap ($^{\circ}$C) & 95\% interval & Newey--West trend ($^{\circ}$C/yr)\\
\midrule
Daily mean (TMK) & +0.455 & 0.432 to 0.477 & $-0.0031$ ($-0.0079$ to 0.0016)\\
Daily minimum (TNK) & +1.120 & 1.086 to 1.153 & +0.0014 ($-0.0077$ to 0.0106)\\
Daily maximum (TXK) & +0.120 & 0.098 to 0.143 & $-0.0029$ ($-0.0069$ to 0.0010)\\
\bottomrule
\end{tabular}
\caption{The gap by variable. Spring has the largest gap for all three.}
\end{table}
\begin{table}[h]
\centering\footnotesize\setlength{\tabcolsep}{3pt}
\begin{tabular}{lccc}
\toprule
After $-$ before (3-year windows) & TMK & TNK & TXK\\
\midrule
Urban move, 2008-07-01 & $-0.433$ ($-0.544$, $-0.292$) & $-0.271$ ($-0.386$, $-0.132$) & $-0.241$ ($-0.370$, $-0.120$)\\
Reference move, 2014-10-22 & $+0.187$ (0.103, 0.267) & $+0.414$ (0.234, 0.595) & $+0.101$ (0.026, 0.173)\\
\bottomrule
\end{tabular}
\caption{Step in the mean gap across each documented relocation, with block-bootstrap 95\% intervals.}
\end{table}
By segment the daily-mean gap is $+0.506$ (1985--2008-06), $+0.337$ (2008-07 to 2014-10) and $+0.417$ (2014-10 onward). The annual
means make the step visible: 0.72, 0.82 and 0.81\,$^{\circ}$C in 2005--2007 against 0.36 and 0.27\,$^{\circ}$C in 2009--2010. An annual
regression with level shifts at 2009 and 2015 and Newey--West errors cannot separate steps from trend with only 40 points (slope
$+0.0019$, 95\% interval $-0.0147$ to $+0.0185$; 2009 step $-0.188$, $-0.485$ to $+0.108$), so the direct window comparison is the
more precise instrument.

\section{Discussion}
The documented 2008 move was described as ``slight'' and within the same district, yet the gap fell by about as much as its
entire full-period mean. The full-period mean therefore averages over three different station histories, and the null
trend result, which a reader could take as ``the urban heat island is not changing'', says little about real change while
instrument and siting steps of this size are present. Two caveats apply: three-year windows also contain real
interannual variability and any true change in the surroundings, so the steps are coincident with, not proven
caused by, the moves; and the annual regression's wide intervals show the data alone cannot homogenise the series.
Physically the pattern is coherent with the urban heat island being mainly a night-time phenomenon, with a gap in the daily
minimum 2.5 times that of the daily mean. A proper estimate of change would need the stations' homogenisation metadata or a
reference network, which this study does not have.

\section{Correction}
The earlier report stated that neither relocation coincides with an unusual jump in the annual-mean-gap series. The annual
series contradicts that. The statement is corrected in the repository (erratum in the report); no frozen registry value changed.

\section{Limitations}
Two stations and one pairing; the reference is an airport, not a rural site; elevation and distance (15.4\,km) differ; only
three temperature variables; the three-year window is an arbitrary choice fixed in the code before the step results were seen; and the results are descriptive, not a homogenisation.

\section*{Reproducibility and AI assistance}
Code, frozen registry, post-hoc extensions and tests:\\
\url{https://github.com/lindgreendavid/climate-twin-frankfurt}\\
Interactive laboratory:\\
\url{https://climate-twin-frankfurt-interactive.lindgreendavid.workers.dev}\\
The analysis code, tests and drafts of this text were produced with substantial assistance from an AI system (Claude, Anthropic);
the author is responsible for the content.

\begin{thebibliography}{9}
\bibitem{newey1987} W.~K. Newey and K.~D. West, A simple, positive semi-definite, heteroskedasticity and autocorrelation
consistent covariance matrix, \emph{Econometrica} \textbf{55}, 703--708 (1987).
\end{thebibliography}
\end{document}
